\section{Related Literature and Research Positioning}
\label{sec:related_work}

\subsection{Airline Recovery Decisions and Multidimensional Operational Consequences}
\label{subsec:rw_recovery}

Airline operations are closely interconnected, and a disruption to an individual flight can affect not only its own schedule execution but also subsequent operations through aircraft turnaround, crew connections, passenger itineraries, and airport resource use. Recovery decisions therefore require coordination across interdependent operational elements rather than treating disruptions solely as deviations from scheduled flight times. \citet{clausenDisruptionManagementAirline2010} review aircraft, crew, and passenger recovery and highlight the interdependencies created by schedules, resource connections, and recovery actions. By incorporating passenger delays and itinerary disruptions into flight recovery, \citet{bratuFlightOperationsRecovery2006} demonstrate why flight-level operating costs alone may not adequately represent the consequences of a recovery plan. \citet{hassanAirlineDisruptionManagement2021} further review developments in integrated recovery and identify continuing gaps between optimization models and airline operating practice. Together, these studies indicate that delay is an important indicator of disruption, but its temporal magnitude captures only part of the wider operational consequences.

This multidimensionality also means that flights with similar delays may have substantially different operational implications. Subsequent aircraft assignments, the number and connections of affected passengers, and resource pressures under different airport and operating conditions can all shape the significance of a disruption. The integrated aircraft--passenger recovery formulation of \citet{arikanIntegratedAircraftPassenger2016} likewise illustrates the relationship between operating adjustments and passenger outcomes. Yet assessing multiple consequences raises a further question: how should outcomes measured in different units, with different managerial meanings and relative importance, be compared? Multiattribute decision analysis provides a basis for making such trade-offs explicit. \citet{keeneyRaiffaDecisionsMultipleObjectives1976} examine preferences and value trade-offs across decision objectives, while \citet{oecdCompositeIndicators2008} emphasize normalization, weighting, and sensitivity analysis in constructing composite indicators. These approaches support transparent multidimensional evaluation, but the weights and aggregation rules still require a defensible managerial interpretation. Standardization alone does not turn a composite measure into an observed economic cost or a decision maker's elicited utility.

Beyond evaluating consequences, airline recovery involves different levels of managerial decision making. During operations, controllers assess disruption importance and current operating conditions to determine which problems deserve attention and whether further adjustments are warranted. \citet{peiDecisionSupportIrregular2021} link real-time flight assessment with subsequent recovery recommendations in a decision-support approach for irregular operations. \citet{jiFlightPrioritiesRecovery2021} incorporate flight priorities and airport-capacity constraints into recovery planning, reflecting trade-offs between operational objectives and constrained resources. Under uncertainty, \citet{zhaoTwoStageAircraftRecovery2023} distinguish the identification of affected flights from subsequent schedule adjustments, illustrating the staged nature of information revelation and recovery planning. Airline recovery is therefore not a single task governed by one uniform comparison, but involves disruption identification, priority assessment, evaluation of recovery opportunities, and operational adjustment.

These tasks are related without having identical decision objects. Consequence severity describes the potential impact of a disruption on the operating system, whereas recoverable value also depends on the adjustment opportunities that remain, the available alternatives, and the effort required to improve outcomes. A highly consequential disruption may warrant early attention without offering much scope for additional local recovery; conversely, a less severe disruption may allow a more readily attainable improvement. The staged-recovery literature provides a basis for recognizing this distinction and, in turn, suggests that different tasks need not rely on operating information in the same way. An accurate assessment of disruption severity does not, by itself, constitute an equally informative assessment of recovery opportunity.

\subsection{Rolling Operational Information and Decision-Based Evaluation}
\label{subsec:rw_uncertainty}

Airline recovery decisions are commonly made before the operating process has finished and while future states remain uncertain. As flights progress, aircraft arrivals, turnaround activities, subsequent flight execution, and airport conditions are continually updated, with previously unresolved events becoming observed facts. \citet{vinkDynamicAircraftRecovery2020} incorporate evolving disruptions and operational conditions into a dynamically updated aircraft-recovery framework, while \citet{leeDynamicDisruptionManagement2020} study recovery decisions under airport operating uncertainty. The treatment of disruption duration and information revelation in \citet{zhaoTwoStageAircraftRecovery2023} likewise underscores that the information available to decision makers changes over time. Operational information in airline recovery thus concerns not only forecasting future states, but also using newly available evidence as operations unfold.

As the operating state evolves, its current description and the history leading to it may offer different bases for assessment. The current state summarizes conditions observable or inferable at a given time, whereas the preceding trajectory may retain information about how delays developed, propagated, and changed. \citet{fleurquinSystemicDelayPropagation2013} study systemic delay propagation through interactions in the airline network, \citet{li2021delayNetwork} characterize delay propagation in the air-traffic network, and \citet{erdemAirlineDelayPropagation2024} examine delay estimation and evolution in daily operations. These studies establish the temporal interconnectedness of airline operations, but do not imply that historical trajectories necessarily add decision value beyond the current state. The usefulness of history depends on how much of the relevant past is already summarized by current information. When the current representation is sufficient for the subsequent assessment, retaining the trajectory may add complexity without a corresponding benefit; when it omits relevant dynamics, history may provide additional insight.

More complete state information, however, does not necessarily translate into better decisions. Statistical prediction is often assessed through estimation error, probability calibration, or uncertainty representation, whereas managerial decisions depend on the consequences of the alternatives selected. \citet{blackwellEquivalentComparisonsExperiments1953} establish a foundational comparison of information structures in statistical decision theory, and \citet{howardInformationValueTheory1966} emphasize that the value of information cannot be separated from the consequences of uncertain events. \citet{avrielWilliamsValueInformation1970} relate information value to stochastic programming. This decision-oriented perspective has also informed more recent predictive modeling: \citet{dontiTaskBasedEndtoEnd2017} incorporate downstream optimization tasks into learning, while \citet{elmachtoubGrigasSmartPredictOptimize2022} show why predictive-error improvements need not translate directly into improved optimization decisions. Despite differences in setting and formulation, these contributions distinguish the quality of information from its performance in a particular decision task.

This distinction is especially pertinent to airline recovery because different tasks require different comparisons. Disruption screening concerns which flights deserve priority relative to other candidates; local recovery concerns whether, and how far, a particular flight should be adjusted given the consequences and effort associated with alternative choices. The former depends on cross-flight comparisons, whereas the latter depends on comparisons among recovery options within a flight. Additional information may change an operating-state assessment without changing either decision, and a changed decision may carry only a small difference in the relevant objective. The managerial value of information consequently depends both on the operating-state distinctions it preserves and on their relevance to the particular decision and evaluation criterion.

\subsection{Decision Objectives and Distributional Information Requirements}
\label{subsec:rw_decision_value}

When future operations remain uncertain, decision support also depends on how that uncertainty is represented. A point prediction offers a compact summary but does not describe the probabilities of alternative realizations. A probability distribution retains this variation, while a joint distribution can additionally represent dependence among uncertain quantities that jointly influence operating consequences. Scenario-based approaches already illustrate the importance of uncertainty in airline operations. \citet{wooScenarioBasedStochastic2021} use stochastic scenarios for flight rescheduling, \citet{ogunsinaRelationalDynamicBayesian2022} model uncertainty and its propagation using a relational dynamic Bayesian network, and \citet{zoutendijkProbabilisticFlightDelay2021} apply probabilistic delay prediction to subsequent airport decisions. In the broader decision-analysis literature, \citet{clemenReillyCorrelationsCopulas1999} explain how marginal distributions and copulas can be used to construct joint probability models, distinguishing uncertainty about individual variables from dependence between them.

A more complete probability representation need not, however, yield a comparable improvement in decision performance. Joint distributions describe the co-occurrence of uncertain states, but their usefulness depends on the role that such joint realizations play in the decision problem. \citet{agrawalPriceCorrelations2012} examine the cost of neglecting correlations in stochastic optimization and show that it depends on problem structure: the resulting loss can be controlled in some classes of problems but can be substantial in others. Thus, the importance of dependence information cannot be inferred solely from the existence of statistical correlation. The amount of structure preserved in a representation and the extent to which that structure improves a decision are different considerations.

The decision objective itself also determines which probability features enter the evaluation of an alternative. For an expected objective that is additive across uncertain variables, preserving the relevant marginal distributions may suffice to evaluate each fixed alternative. By contrast, when consequences depend on joint realizations or the effect of an adjustment is constrained by relationships among uncertain states, dependence can affect the expected performance of different alternatives. \citet{clemenReillyCorrelationsCopulas1999} examine dependence modeling and sensitivity to correlations in decision analysis, while \citet{wangDyerCopulasDecisionTrees2012} develop a copula-based approach to representing dependent uncertainties in multivariate decision trees. Together with the problem-dependent costs of ignoring correlation studied by \citet{agrawalPriceCorrelations2012}, these contributions indicate why the relevance of dependence must be assessed in relation to how an objective translates uncertain states into consequences. Different managerial objectives can therefore place different demands on the same underlying probability representation.

Even when dependence changes the evaluation of possible outcomes, it need not change the final choice. Decisions involve comparisons among feasible alternatives, so what matters for selection is whether a change in probability representation alters their relative performance. A representation can change expected objective values while preserving the preferred alternative; conversely, even a modest relative change can reverse the choice when alternatives are closely matched. The decision-consequence perspective of \citet{howardInformationValueTheory1966}, together with the downstream objective-based evaluation developed by \citet{elmachtoubGrigasSmartPredictOptimize2022}, provides a basis for distinguishing informational refinement from decision improvement. The managerial value of distributional information must therefore be considered in relation to probability structure, decision objectives, and the comparisons among alternatives, rather than judged by representational completeness alone.

\subsection{Research Gap}
\label{subsec:rw_gap}

Existing research provides important foundations for uncertainty modeling in airline recovery and for decision-based information evaluation. The connection between these perspectives in specific airline operating decisions, however, warrants further empirical examination.

For the same operating process, there remains a need for systematic evidence on when simpler operating signals or probability representations are sufficient for different recovery tasks. In particular, it is not yet well established under what conditions a delay-based signal preserves the recovery priorities supported by multidimensional consequences, when current-state information or marginal uncertainty can substitute for historical trajectories or joint dependence, and how the value of these distinctions varies with managerial objectives and operating conditions. The central issue is not merely whether a richer representation improves the description of the operating state, but whether the distinctions it preserves change the comparison among recovery alternatives and the consequences of the resulting decisions.

Against this background, the present study examines the roles of operating information and its representations in recovery screening and local-recovery decisions under rolling airline operations. It further considers how decision objectives and operating conditions shape their decision relevance, with the aim of identifying when simpler information supports the same decisions and when richer operating information offers incremental value.
